\documentclass[11pt]{article}
\usepackage{amsmath, amssymb, amsthm}
\usepackage[margin=1in]{geometry}
\usepackage{algorithm}
\usepackage{algpseudocode}
\usepackage{booktabs}

\newtheorem{theorem}{Theorem}
\newtheorem{lemma}[theorem]{Lemma}

\newcommand{\LC}{\mathrm{LC}}

\title{Problem 412: Gnomon Numbering}
\author{}
\date{}

\begin{document}
\maketitle

\section{Problem Statement}

For integers $m$ and $n$ with $0 \le n < m$, define $L(m,n)$ as the gnomon formed by an $m \times m$ grid with the top-right $n \times n$ portion removed. $\LC(m,n)$ counts valid fillings with $1, 2, \ldots, m^2 - n^2$ such that each cell is smaller than the cell below and to its right.

\textbf{Find:} $\LC(10000, 5000) \bmod 76543217$.

\section{Mathematical Foundation}

\begin{theorem}[Hook Length Formula, Frame--Robinson--Thrall 1954]
For a partition $\lambda$ with $N$ cells, the number of standard Young tableaux of shape~$\lambda$ is
\[
f^\lambda = \frac{N!}{\displaystyle\prod_{(i,j) \in \lambda} h(i,j)}
\]
where $h(i,j) = \mathrm{arm}(i,j) + \mathrm{leg}(i,j) + 1$.
\end{theorem}

\begin{proof}
This is a classical result in algebraic combinatorics; see Frame, Robinson, and Thrall (1954). The proof proceeds via a bijection between standard Young tableaux and hook-removal sequences, or equivalently through the representation theory of the symmetric group.
\end{proof}

\begin{lemma}[Gnomon as Young Diagram]
The gnomon $L(m,n)$ corresponds to the partition $\lambda = (m^{m-n}, (m-n)^n)$, and a valid filling is a standard Young tableau of shape~$\lambda$.
\end{lemma}

\begin{proof}
Setting $a = m - n$, the first $a$ rows have length $m$ and the last $n$ rows have length $a$. The filling constraint (each cell smaller than neighbors below and to the right) is the standard Young tableau condition.
\end{proof}

\begin{lemma}[Hook Length Decomposition]
With $a = m - n$, the cells decompose into three rectangular regions. The hook length at $(i,j)$ is:
\begin{center}
\begin{tabular}{@{}llll@{}}
\toprule
Region & Rows & Columns & $h(i,j)$ \\
\midrule
1 & $i \in [0,a)$ & $j \in [0,a)$ & $2m - i - j - 1$ \\
2 & $i \in [0,a)$ & $j \in [a,m)$ & $m + a - i - j - 1$ \\
3 & $i \in [a,m)$ & $j \in [0,a)$ & $a + m - i - j - 1$ \\
\bottomrule
\end{tabular}
\end{center}
\end{lemma}

\begin{proof}
In Region~1: $\mathrm{arm}(i,j) = m - j - 1$ and $\mathrm{leg}(i,j) = m - i - 1$, giving $h = 2m - i - j - 1$. In Region~2: $\mathrm{arm} = m - j - 1$ and $\mathrm{leg} = a - i - 1$, so $h = m + a - i - j - 1$. In Region~3: $\mathrm{arm} = a - j - 1$ and $\mathrm{leg} = m - i - 1$, so $h = a + m - i - j - 1$.
\end{proof}

\begin{theorem}[Row-wise Factorization]
The product of all hook lengths factorizes as:
\[
\prod_{(i,j) \in \lambda} h(i,j) =
\prod_{i=0}^{a-1} \frac{(2m{-}i{-}1)!}{(m{+}n{-}i{-}1)!}
\cdot \prod_{i=0}^{a-1} \frac{(m{-}i{-}1)!}{(a{-}i{-}1)!}
\cdot \prod_{i=a}^{m-1} \frac{(a{+}m{-}i{-}1)!}{(m{-}i{-}1)!}
\]
\end{theorem}

\begin{proof}
For row $i$ in Region~1, $\prod_{j=0}^{a-1}(2m-i-j-1) = (2m-i-1)!/(2m-i-a-1)! = (2m-i-1)!/(m+n-i-1)!$. The other two regions follow by identical telescoping of consecutive-integer products.
\end{proof}

\section{Algorithm}

\begin{algorithm}[H]
\caption{Compute $\LC(m, n) \bmod p$}
\begin{algorithmic}[1]
\State $a \gets m - n$, \quad $N \gets m^2 - n^2$
\State Precompute $k! \bmod p$ and $(k!)^{-1} \bmod p$ for $k = 0, 1, \ldots, N$
\State $D \gets 1$
\For{$i = 0$ to $a - 1$} \Comment{Region 1}
    \State $D \gets D \cdot (2m{-}i{-}1)! \cdot ((m{+}n{-}i{-}1)!)^{-1} \pmod{p}$
\EndFor
\For{$i = 0$ to $a - 1$} \Comment{Region 2}
    \State $D \gets D \cdot (m{-}i{-}1)! \cdot ((a{-}i{-}1)!)^{-1} \pmod{p}$
\EndFor
\For{$i = a$ to $m - 1$} \Comment{Region 3}
    \State $D \gets D \cdot (a{+}m{-}i{-}1)! \cdot ((m{-}i{-}1)!)^{-1} \pmod{p}$
\EndFor
\State \Return $N! \cdot D^{-1} \pmod{p}$
\end{algorithmic}
\end{algorithm}

\section{Complexity Analysis}

\begin{itemize}
  \item \textbf{Time:} $O(N) = O(m^2)$ for factorial precomputation; $O(m)$ for the three loops. Total: $O(m^2)$.
  \item \textbf{Space:} $O(N) = O(m^2)$ for factorial tables.
\end{itemize}

\section{Answer}

\[
\boxed{38788800}
\]

\end{document}
